% Minimal paper skeleton for the paper-writer agent.
%
% Copy the venue style file you need next to this file (see latex/<venue>/),
% then uncomment the matching \usepackage line below.
%
% The agent fills in each section to what it must ESTABLISH (see
% agents/claude/paper-writer.md), not to a page budget.

\documentclass{article}

% --- pick one venue style ---
\usepackage[preprint]{neurips_2025}  % NeurIPS  -> latex/neurips/neurips_2025.sty
%\usepackage{aaai2027}                % AAAI     -> latex/aaai/aaai2027.sty (+ .bst)
%\usepackage{acl}                    % *ACL     -> latex/acl/acl.sty (+ acl_natbib.bst)

\usepackage{amsmath,amssymb}
\usepackage{booktabs}
\usepackage{hyperref}

\title{Paper Title}
\author{Author Name}

\begin{document}
\maketitle

\begin{abstract}
Write this LAST. Four to six sentences: a miniature of the whole paper.
Opening is conditional -- familiar problem + surprising result leads with the
achievement; unfamiliar framing opens with something every reader agrees with.
Then why it is hard, what you did (one teaser sentence), the evidence with
your most remarkable number. State assumptions and scope honestly.
\end{abstract}

\section{Introduction}
% Establishes: the problem is important, the problem is hard, and you solved it.
% Write the contributions FIRST -- two to four refutable bullets, each carrying
% its section pointer. Every claim here must be discharged by a named section
% (the claim ledger); check each one before finishing.

\section{Background / Related Work}
% Establishes the DIFF, not the inventory. Group by methodological commitment,
% never paper-by-paper. Be generous with credit. Cite only what you have opened.

\section{Method}
% Establishes self-sufficiency: output and key idea FIRST, then mechanism.
% Intuition and a worked toy example before the general case. Design
% alternatives go to the appendix, not inline.

\section{Experimental Setup}
% Establishes reproducibility: datasets and splits, task and metrics, baselines
% AND how they were tuned, hyperparameters and how selected, compute, seeds and
% number of runs. Always state whether a number is a single run, a max, or a mean.

\section{Results}
% Establishes that the claim rests on data and logic. One subsection per claim,
% each opening with a signpost. Narrate, do not re-key the tables. Report
% negative and mixed results plainly.

\section{Analysis and Limitations}
% Write this BEFORE the abstract. Establishes what the results MEAN and where
% they stop being true. Limitations name the MECHANISM that bounds the method
% ("relies on limited morphology, as in English"), not surface facts about the
% experiment grid ("we only evaluated on English").

\section{Conclusion}
% Establishes what changed, compressed and sharpened. Never merely restate the
% abstract. Future work framed as a limitation, never a wish list. No puffery.

\bibliographystyle{plain}
\bibliography{references}  % provide references.bib, or swap style per venue

\end{document}
